\documentclass[../DoAn.tex]{subfiles}
\usepackage{multirow}
\usepackage{longtable}
\usepackage{placeins}
\usepackage{array}
\usepackage{graphicx}
\usepackage{float}
\usepackage{caption}
\usepackage{tabularx}
\usepackage{adjustbox}
\usepackage[hyphens]{url}
\setlist[itemize,2]{label=\textbullet}
\begin{document}
% Chương này có độ dài từ 9 đến 11 trang.
% Với phương pháp phân tích thiết kế hướng đối tượng, sinh viên sử dụng biểu đồ use case theo hướng dẫn của template này. Với các phương pháp khác, sinh viên trao đổi với giáo viên hướng dẫn để đổi tên và sắp xếp lại đề mục cho phù hợp. Ví dụ, thay vì sử dụng biểu đồ use case, sinh viên đi theo hướng tiếp cận Agile có thể dùng User Story.
% \textbf{Lưu ý}: Mỗi chương nên có thêm 1 đoạn mở đầu chương và kết thúc chương, mở đầu giới thiệu những nội dung sẽ trình bày trong chương, kết thúc tổng kết lại các nội dung đã trình bày

Trong chương này, em sẽ thực hiện khảo sát các hệ thống thương mại điện tử đã có trên thị trường, tìm ra ưu điểm và hạn chế còn tồn đọng của các hệ thống đã có, từ đó phân tích và đưa ra các chức năng chính và các yêu cầu phi chức năng của hệ thống

\section{Khảo sát hiện trạng}
\label{section:2.1}
% Thông thường, khảo sát chi tiết về hiện trạng và yêu cầu của phần mềm sẽ được lấy từ ba nguồn chính, đó là (i) người dùng/khách hàng, (ii) các hệ thống đã có, (iii) và các ứng dụng tương tự.
% Sinh viên cần tiến hành phân tích, so sánh, đánh giá chi tiết ưu nhược điểm của các sản phẩm/nghiên cứu hiện có. Sinh viên có thể lập bảng so sánh nếu cần thiết. Kết hợp với khảo sát người dùng/khách hàng (nếu có), sinh viên nêu và mô tả sơ lược các tính năng phần mềm quan trọng cần phát triển.


\subsection{Khảo sát hệ thống hiện có}
\label{subsection:2.1.1}
Hai ứng dụng thương mại điện tử nổi bật trên nền tảng di động hiện nay là \textbf{Shopee} và \textbf{TikTok Shop}. Đây là những nền tảng lớn, có lượng người dùng đông đảo và nhiều chức năng phong phú, là cơ sở quan trọng để so sánh, phân tích và đề xuất cho hệ thống cần phát triển.

\subsubsection*{Shopee}
\label{subsection:2.1.1.1}
\begin{itemize}
  \item \textbf{Ưu điểm:}
  \begin{itemize}
    \item Giao diện trực quan, dễ sử dụng cho cả người bán và người mua.
    \item Hệ thống phân loại sản phẩm rõ ràng, nhiều bộ lọc thông minh.
    \item Tích hợp hệ thống thanh toán nội bộ (ShopeePay), ví điện tử và liên kết ngân hàng.
    \item Chức năng Flash Sale, Voucher hấp dẫn.
    \item Tính năng livestream bán hàng hỗ trợ người bán tiếp cận người dùng tốt hơn.
    \item Chính sách bảo vệ người mua, hoàn tiền rõ ràng.
  \end{itemize}

  \item \textbf{Nhược điểm:}
  \begin{itemize}
    \item Giao diện đôi khi quá nhiều thông tin gây rối mắt.
    \item Tốc độ tải chậm ở một số thiết bị cấu hình yếu.
  \end{itemize}
\end{itemize}

\subsubsection*{TikTok Shop}
\label{subsection:2.1.1.2}
\begin{itemize}
  \item \textbf{Ưu điểm:}
  \begin{itemize}
    \item Tích hợp trực tiếp với video ngắn giúp tăng trải nghiệm mua hàng giải trí.
    \item Khả năng tiếp cận sản phẩm thông qua KOL/influencer hiệu quả.
    \item Hệ thống gợi ý sản phẩm dựa trên AI theo hành vi người dùng.
    \item Liên kết bán hàng – video – thanh toán mượt mà trên cùng một nền tảng.
  \end{itemize}

  \item \textbf{Nhược điểm:}
  \begin{itemize}
    \item Giao diện quản lý đơn hàng và sản phẩm của người bán còn thiếu chuyên nghiệp.
    \item Hệ thống phân loại và tìm kiếm sản phẩm chưa tối ưu.
  \end{itemize}
\end{itemize}

\subsection{Các chức năng cần có cho hệ thống}
\label{subsection:2.1.2}
Dựa trên phân tích trên, hệ thống thương mại điện tử cần xây dựng nên có các chức năng cốt lõi sau:

\begin{itemize}
  \item Đăng ký/đăng nhập, quản lý, cài đặt tài khoản
  \item Tìm kiếm, lọc, xem chi tiết sản phẩm.
  \item Giỏ hàng thanh toán đa phương thức (ví điện tử, ngân hàng, COD).
  \item Quản lý đơn hàng, trạng thái vận chuyển.
  \item Đánh giá, bình luận sản phẩm.
  \item Tạo và quản lý voucher, chương trình khuyến mãi.
  \item Quản trị viên quản lý người dùng, đơn hàng, sản phẩm.
  \item Tích hợp bên thứ ba: API thanh toán (Stripe, Momo), API giao hàng (GHN, GHTK).
\end{itemize}

\subsection{Tính năng đề xuất}
\label{subsection:2.1.3}
Ngoài các tính năng chính cốt lõi trên, hệ thống của đồ án sẽ có thêm một số \textbf{tính năng} sau:

\begin{itemize}
  \item \textbf{Gợi ý sản phẩm cá nhân hóa:} Sử dụng mô hình học máy để đưa ra các gợi ý cá nhân hóa cho mỗi người dùng
  \item \textbf{Chatbot tự động hóa quy trình mua sắm:} Chatbot tự động hóa quy trình mua sắm khi sẽ gợi ý cho người dùng những hành động người dùng chỉ cần nhập vài câu lệnh đơn giản hoặc click quick button là có thể tìm kiếm, xem sản phẩm, thêm vào giỏ hàng, xem giỏ hàng hay đặt hàng ngay lập tức, vô cùng nhanh chóng và tiện lợi.
  \item \textbf{Tạo nhóm săn sale:} Cho phép người dùng có thể kết bạn, tạo nhóm và mời bạn bè vào nhóm để cùng nhau săn sale và những voucher hấp dẫn của hệ thống
  \item \textbf{Tham gia các nhiệm vụ:} Hệ thống tạo các nhiệm vụ có thời gian hoàn thành đếm ngược và độ khó khác nhau khuyến khích người dùng tham gia trải nghiệm tương tác ứng dụng nhiều hơn để nhận các phần quà hấp dẫn.
\end{itemize}



\section{Tổng quan chức năng}
\label{section:2.2}
% Phần \ref{section:2.2} này có nhiệm vụ tóm tắt các chức năng của phần mềm. Trong phần này, sinh viên lưu ý chỉ mô tả chức năng mức cao (tổng quan) mà không đặc tả chi tiết cho từng chức năng. Đặc tả chi tiết được trình bày trong phần \ref{section:2.3}.


% Sinh viên vẽ biểu đồ use case tổng quan và giải thích các tác nhân tham gia là gì, nêu vai trò của từng tác nhân, và mô tả ngắn gọn các use case chính.


\newpage
\subsection{Biểu đồ use case tổng quát}
\label{subsection:2.2.1}

\begin{figure}[!htbp]
\centering
\includegraphics[width=1.2\linewidth]{Hinhve/uc.jpg}
\caption{Biểu đồ use case tổng quan}
\label{fig:usecase}
\end{figure}
\newpage
Mô tả tác nhân: Hệ thống gồm 4 tác nhân là: Người dùng, Quản trị viên, API thanh toán (Stripe) và API giao hàng (GHN)


\begin{table}[h!]
\centering
\begin{tabular}{|c|c|p{9cm}|}
\hline
\textbf{STT} & \textbf{Tên tác nhân} & \multicolumn{1}{c|}{\textbf{Mô tả}} \\ \hline
\textbf{1} & Người dùng &
Là khách hàng sau khi đã đăng ký tài khoản và đăng nhập thành công. Người dùng có thể trải nghiệm đầy đủ các tính năng hấp dẫn của hệ thống thương mại điện tử bao gồm: xem và tìm kiểm sản phẩm, xem sản phẩm được cá nhân hóa, xem cửa hàng, category, thương hiệu nổi bật hiện có trên sàn, thêm vào giỏ hàng, đặt hàng, thanh toán, nhận thông báo trong hệ thống, xem và viết bình luận về sản phẩm, thu thập và sử dụng voucher, tham gia nhiệm vụ để tích điểm thưởng và đổi điểm thưởng lấy quà hấp dẫn. Ngoài ra người dùng có thể xem và theo dõi đơn hàng, mua lại đơn hàng, quản lý bạn bè, quản lý địa chỉ, quản lý cài đặt tài khoản, quản lý sản phẩm yêu thích. Và đặc biệt là tham gia hoạt động tạo nhóm săn sale cùng bạn bè vô cùng hấp dẫn và mang lại sự thích thú và tối ưu trong quá trình mua sắm.   \\ \hline

\textbf{2} & Quản trị viên hệ thống &
Có thể đăng nhập vào web Admin để thực hiện các công việc quản lý, xem thống kê sản phẩm, đơn hàng, người dùng, shops, thương hiệu, danh mục, voucher, nhiệm vụ trên hệ thống. \\ \hline

\textbf{3} & API thanh toán (Stripe) &
Cung cấp chức năng thanh toán điện tử cho người dùng. Hệ thống sẽ gọi API này để xử lý thanh toán đơn hàng một cách an toàn, nhanh chóng và bảo mật. \\ \hline

\textbf{4} & API giao hàng (GHN) &
Cung cấp dịch vụ giao hàng. Hệ thống gọi API này để tự động tạo đơn vận chuyển và lấy thông tin vận chuyển liên quan đến đơn hàng. \\ \hline
\end{tabular}
\caption{Mô tả các tác nhân trong hệ thống}
\label{tab:actors}
\end{table}



\FloatBarrier



Hình \ref{fig:usecase} trình bày biểu đồ ca sử dụng tổng quan của hệ thống đối với các tác nhân đã được mô tả vắn tắt các ca sử dụng ở Bảng \ref{tab:actors}
\newpage
\subsection{Biểu đồ use case phân rã "Quản lý đơn hàng"}
\label{subsection:2.2.2}
Hình \ref{fig:donhang} mô tả phân rã của use case "Quản lý đơn hàng". Với use case này, tác nhân là quản trị viên của hệ thống, có thể xem toàn bộ các đơn hàng của toàn bộ người dùng trong hệ thống, cập nhật trạng thái của đơn hàng, tìm kiếm và lọc đơn hàng, xem thống kê tổng doanh thu của hệ thống.
\begin{figure}[h!]
\centering
\includegraphics[width=0.75\linewidth]{Hinhve/Quản lý đơn hàng.jpg}
\caption{Biểu đồ use case phân rã "Quản lý đơn hàng"}
\label{fig:donhang}
\end{figure}
\FloatBarrier
\subsection{Biểu đồ use case phân rã "Quản lý người dùng"}
\label{subsection:2.2.3}
Hình \ref{fig:nguoidung} mô tả phân rã của use case "Quản lý người dùng". Với use case này, tác nhân là quản trị viên của hệ thống, có thể xem toàn bộ thông tin của tất cả người dùng trong hệ thống, cập nhật thông tin người dùng, cập nhật role của họ, tìm kiếm người dùng. Đặc biệt quản trị viên có thể khóa người dùng và mở khóa người dùng nếu phát hiện các hành vi gian lận mà không cần xóa tài khoản đó, có thể dễ dàng mở khóa lại người dùng sau khi hợp tác hoặc khi vấn đề đã được giải quyết.
\begin{figure}[h!]
\centering
\includegraphics[width=0.75\linewidth]{Hinhve/datn/Quản lý người dùng.jpg}
\caption{Biểu đồ use case phân rã "Quản lý người dùng"}
\label{fig:nguoidung}
\end{figure}
\newpage
\FloatBarrier
\subsection{Biểu đồ use case phân rã "Quản lý cửa hàng"}
\label{subsection:2.2.4}
Hình \ref{fig:cuahang} mô tả phân rã của use case "Quản lý cửa hàng". Với use case này, tác nhân là quản trị viên của hệ thống, có thể xem toàn bộ thông tin của tất cả các cửa hàng đã đăng ký trong hệ thống, tìm kiếm cửa hàng, thêm 1 cửa hàng mới, cập nhật thông tin cửa hàng, xóa cửa hàng và xem tổng số cửa hàng trong hệ thống.
\begin{figure}[h!]
\centering
\includegraphics[width=0.75\linewidth]{Hinhve/datn/Quản lý cửa hàng.jpg}
\caption{Biểu đồ use case phân rã "Quản lý cửa hàng"}
\label{fig:cuahang}
\end{figure}
\newpage
\FloatBarrier
\subsection{Biểu đồ use case phân rã "Quản lý category"}
\label{subsection:2.2.5}
Hình \ref{fig:category} mô tả phân rã của use case "Quản lý category". Với use case này, tác nhân là quản trị viên của hệ thống, có thể xem toàn bộ thông tin của tất cả các category đang có trong hệ thống, tìm kiếm category, thêm 1 category mới, cập nhật thông tin category đã có, xóa category và xem tổng số category trong hệ thống.
\begin{figure}[h!]
\centering
\includegraphics[width=0.75\linewidth]{Hinhve/datn/Quản lý category.jpg}
\caption{Biểu đồ use case phân rã "Quản lý category"}
\label{fig:category}
\end{figure}
\newpage
\FloatBarrier
\subsection{Biểu đồ use case phân rã "Quản lý nhiệm vụ"}
\label{subsection:2.2.6}
Hình \ref{fig:nhiemvu} mô tả phân rã của use case "Quản lý nhiệm vụ". Với use case này, tác nhân là quản trị viên của hệ thống, có thể xem toàn bộ thông tin của tất cả các nhiệm vụ đang có trong hệ thống, tìm kiếm nhiệm vụ, lọc nhiệm vụ theo 1 type cụ thể, thêm 1 nhiệm vụ mới, cập nhật thông tin nhiệm vụ đã có, xóa nhiệm vụ và xem tổng số nhiệm vụ đã được tạo trong hệ thống.
\begin{figure}[h!]
\centering
\includegraphics[width=0.75\linewidth]{Hinhve/datn/Quản lý nhiệm vụ.jpg}
\caption{Biểu đồ use case phân rã "Quản lý nhiệm vụ"}
\label{fig:nhiemvu}
\end{figure}
\FloatBarrier
\subsection{Biểu đồ use case phân rã "Xem thống kê"}
\label{subsection:2.2.7}
Hình \ref{fig:xemthongke} mô tả phân rã của use case "Xem thống kê". Với use case này, tác nhân là quản trị viên của hệ thống, có thể xem tất cả các biểu đồ thống kê trong hệ thống. Hệ thống có 2 loại biểu đồ chính là biểu đồ tròn: gồm các biểu đồ thống kê đơn hàng theo status, thống kê nhiệm vụ theo type, thống kê voucher theo type, thống kê tỉ lệ sản phẩm theo category hoặc brand hoặc shop. Loại thứ hai là biểu đồ đường, diễn tả sự biến động theo thời gian gồm: thống kê đơn hàng, thống kê nhiệm vụ, thống kê voucher theo thời gian.
\begin{figure}[h!]
\centering
\includegraphics[width=0.75\linewidth]{Hinhve/datn/Xem thống kê.jpg}
\caption{Biểu đồ use case phân rã "Xem thống kê"}
\label{fig:xemthongke}
\end{figure}
\newpage
\FloatBarrier
\subsection{Biểu đồ use case phân rã "Quản lý voucher"}
\label{subsection:2.2.8}
Hình \ref{fig:voucher} mô tả phân rã của use case "Quản lý voucher". Với use case này, tác nhân là quản trị viên của hệ thống, có thể xem toàn bộ thông tin của tất cả các voucher đang có trong hệ thống, tìm kiếm voucher, lọc voucher theo 1 type cụ thể, thêm 1 voucher mới, cập nhật thông tin voucher đã có, xóa voucher và xem tổng số voucher đã được tạo trong hệ thống.
\begin{figure}[h!]
\centering
\includegraphics[width=0.75\linewidth]{Hinhve/datn/Quản lý voucher.jpg}
\caption{Biểu đồ use case phân rã "Quản lý voucher"}
\label{fig:voucher}
\end{figure}
\FloatBarrier
\subsection{Biểu đồ use case phân rã "Quản lý sản phẩm"}
\label{subsection:2.2.9}
Hình \ref{fig:quanlysanpham} mô tả phân rã của use case "Quản lý sản phẩm". Với use case này, tác nhân là quản trị viên của hệ thống, có thể xem toàn bộ thông tin của tất cả các sản phẩm đang có trong hệ thống, tìm kiếm sản phẩm, thêm 1 sản phẩm mới, cập nhật thông tin sản phẩm đã có, xóa sản phẩm và xem tổng số sản phẩm đang có trong hệ thống.
\begin{figure}[h!]
\centering
\includegraphics[width=0.75\linewidth]{Hinhve/datn/Quản lý sản phẩm.jpg}
\caption{Biểu đồ use case phân rã "Quản lý sản phẩm"}
\label{fig:quanlysanpham}
\end{figure}
\newpage
\FloatBarrier
\subsection{Biểu đồ use case phân rã "Quản lý thương hiệu"}
\label{subsection:2.2.10}
Hình \ref{fig:quanlythuonghieu} mô tả phân rã của use case "Quản lý thương hiệu". Với use case này, tác nhân là quản trị viên của hệ thống, có thể xem toàn bộ thông tin của tất cả các thương hiệu đang có trong hệ thống, tìm kiếm thương hiệu, thêm 1 thương hiệu mới, cập nhật thông tin thương hiệu đã có, xóa thương hiệu và xem tổng số thương hiệu đang có trong hệ thống.
\begin{figure}[h!]
\centering
\includegraphics[width=0.75\linewidth]{Hinhve/datn/Quản lý thương hiệu.jpg}
\caption{Biểu đồ use case phân rã "Quản lý thương hiệu"}
\label{fig:quanlythuonghieu}
\end{figure}
\FloatBarrier
\subsection{Biểu đồ use case phân rã "Quản lý bạn bè"}
\label{subsection:2.2.11}
Hình \ref{fig:quanlybanbe} mô tả phân rã của use case "Quản lý bạn bè". Với use case này, tác nhân là người dùng hệ thống, có thể xem danh sách bạn bè hiện tại, gửi yêu cầu kết bạn tới những người khác cũng đang hoạt động trong hệ thống, chấp nhận yêu cầu kết bạn từ người khác gửi đến và hủy kết bạn với một người nào đó.
\begin{figure}[h!]
\centering
\includegraphics[width=0.75\linewidth]{Hinhve/datn/Quản lý bạn bè.jpg}
\caption{Biểu đồ use case phân rã "Quản lý bạn bè"}
\label{fig:quanlybanbe}
\end{figure}
\newpage
\FloatBarrier
\subsection{Biểu đồ use case phân rã "Quản lý giỏ hàng"}
\label{subsection:2.2.12}
Hình \ref{fig:quanlygiohang} mô tả phân rã của use case "Quản lý giỏ hàng". Với use case này, tác nhân là người dùng hệ thống, có thể xem toàn bộ thông tin của tất cả các sản phẩm đang có trong giỏ hàng, thêm 1 sản phẩm mới vào giỏ hàng, cập nhật số lượng của các sản phẩm trong giỏ hàng, xóa sản phẩm khỏi giỏ hàng và xem giá trị của giỏ hàng hiện tại.
\begin{figure}[h!]
\centering
\includegraphics[width=0.75\linewidth]{Hinhve/datn/Quản lý giỏ hàng.jpg}
\caption{Biểu đồ use case phân rã "Quản lý giỏ hàng"}
\label{fig:quanlygiohang}
\end{figure}
\FloatBarrier
\subsection{Biểu đồ use case phân rã "Quản lý địa chỉ"}
\label{subsection:2.2.13}
Hình \ref{fig:quanlydiachi} mô tả phân rã của use case "Quản lý địa chỉ". Với use case này, tác nhân là người dùng hệ thống, có thể xem toàn bộ thông tin của tất cả các địa chỉ mà họ đã tạo, thêm 1 địa chỉ mới, cập nhật thông tin của 1 địa chỉ cũ, xóa 1 địa chỉ và chọn một địa chỉ làm mặc định khi thanh toán đơn hàng
\begin{figure}[h!]
\centering
\includegraphics[width=0.75\linewidth]{Hinhve/datn/Quản lý địa chỉ.jpg}
\caption{Biểu đồ use case phân rã "Quản lý địa chỉ"}
\label{fig:quanlydiachi}
\end{figure}
\newpage
\FloatBarrier
\subsection{Biểu đồ use case phân rã "Tham gia nhóm săn sale"}
\label{subsection:2.2.14}
Hình \ref{fig:thamgianhomsansale} mô tả phân rã của use case "Tham gia nhóm săn sale". Với use case này, tác nhân là người dùng hệ thống, có thể xem toàn bộ thông tin của các nhóm săn sale mà họ đã tạo hoặc họ đã tham gia từ lời mời của bạn bè, tạo 1 nhóm săn sale mới, mời bạn bè tham gia nhóm săn sale mà mình vừa tạo, chấp nhận lời mời để tham gia nhóm săn sale cùng bạn bè, và xóa nhóm săn sale mà mình đã tạo.
\begin{figure}[h!]
\centering
\includegraphics[width=0.75\linewidth]{Hinhve/datn/Tham gia nhóm săn sale.jpg}
\caption{Biểu đồ use case phân rã "Tham gia nhóm săn sale"}
\label{fig:thamgianhomsansale}
\end{figure}
\FloatBarrier
\subsection{Biểu đồ use case phân rã "Quản lý sản phẩm yêu thích"}
\label{subsection:2.2.15}
Hình \ref{fig:quanlysanphamyeuthich} mô tả phân rã của use case "Quản lý sản phẩm yêu thích". Với use case này, tác nhân là người dùng hệ thống, có thể xem toàn bộ thông tin của các sản phẩm mà họ đã thêm vào danh sách yêu thích, xóa một sản phẩm khỏi danh sách yêu thích hiện tại và thêm 1 sản phẩm mới vào danh sách yêu thích.
\begin{figure}[h!]
\centering
\includegraphics[width=0.75\linewidth]{Hinhve/datn/Quản lý sản phẩm yêu thích.jpg}
\caption{Biểu đồ use case phân rã "Quản lý sản phẩm yêu thích"}
\label{fig:quanlysanphamyeuthich}
\end{figure}
\newpage
\FloatBarrier
\subsection{Biểu đồ use case phân rã "Quản lý tài khoản}
\label{subsection:2.2.16}
Hình \ref{fig:quanlytaikhoan} mô tả phân rã của use case "Quản lý tài khoản". Với use case này, tác nhân là người dùng hệ thống, có thể xem thông tin cá nhân của bản thân, bổ sung hoặc cập nhật thông tin cá nhân, chọn ngôn ngữ sẽ hiển thị trong ứng dụng: tiếng Anh hoặc tiếng Việt, chọn chế độ hiển thị trong ứng dụng là chế độ sáng(light mode) hoặc chế độ tối (dark mode) thì giao diện của ứng dụng sẽ cập nhật tương ứng.
\begin{figure}[h!]
\centering
\includegraphics[width=0.75\linewidth]{Hinhve/datn/Quản lý tài khoản.jpg}
\caption{Biểu đồ use case phân rã "Quản lý tài khoản"}
\label{fig:quanlytaikhoan}
\end{figure}

\clearpage
% Nếu sản phẩm/hệ thống cần xây dựng có quy trình nghiệp vụ quan trọng/đáng chú ý, sinh viên cần mô tả và vẽ biểu đồ hoạt động minh họa quy trình nghiệp vụ đó. Sinh viên lưu ý đây không phải là luồng sự kiện của từng use case, mà là luồng hoạt động kết hợp nhiều use case để thực hiện một nghiệp vụ nào đó.

% Ví dụ, một hệ thống quản lý thư viện có quy trình nghiệp vụ mượn trả với mô tả sơ bộ như sau: Sinh viên làm thẻ mượn, sau đó sinh viên đăng ký mượn sách, thủ thư cho mượn, và cuối cùng sinh viên trả lại sách cho thư viện. Một hệ thống có thể có một vài quy trình nghiệp vụ quan trọng như vậy.


% Với mỗi use case mức cao trong biểu đồ use case tổng quan, sinh viên tạo một mục riêng như mục \ref{subsection:2.2.2} và tiến hành phân rã use case đó. Lưu ý tên use case cần phân rã trong biểu đồ use case tổng quan phải khớp với tên đề mục.

% Trong mỗi mục như vậy, sinh viên vẽ và giải thích ngắn gọn các use case phân rã.


\section{Đặc tả chức năng}
\label{section:2.3}
% Sinh viên lựa chọn từ 4 đến 7 use case quan trọng nhất của đồ án để đặc tả chi tiết. Mỗi đặc tả bao gồm ít nhất các thông tin sau: (i) Tên use case, (ii) Luồng sự kiện (chính và phát sinh), (iii) Tiền điều kiện, và (iv) Hậu điều kiện. Sinh viên chỉ vẽ bổ sung biểu đồ hoạt động khi đặc tả use case phức tạp.
\subsection{Đặc tả use case "Xem sản phẩm"}
\hfill

\begin{table}[ht]
\centering
\begin{tabular}{|llrl|}
\hline
\multicolumn{1}{|l|}{\textbf{Mã Use case}}    & \multicolumn{1}{l|}{1} & \multicolumn{1}{r|}{\textbf{Tên Use case}} &  \multicolumn{1}{r|}{Xem sản phẩm } \\ \hline

\multicolumn{1}{|l|}{\textbf{Tác nhân}}       & \multicolumn{3}{l|}{Người dùng}                             \\ \hline
\multicolumn{1}{|l|}{\textbf{Tiền điều kiện}} & \multicolumn{3}{l|}{Người dùng đã đăng nhập vào ứng dụng}                                \\ \hline
\multicolumn{4}{|l|}{}
                    \\ 
\multicolumn{4}{|l|}{
\begin{tabular}{|c|l|c|l|}
\hline
\multirow{6}{*}{\textbf{Luồng sự kiện chính}}                                                   & \textbf{STT} & \begin{tabular}[c]{@{}c@{}}\textbf{Thực hiện} \\ \textbf{bởi}\end{tabular} & \textbf{Hành động}                                                                                                 \\ \cline{2-4} 
                                                                                       & 1   & Người dùng    & Truy cập danh sách sản phẩm                                                        \\ \cline{2-4} 
                                                                                       & 2   & Hệ thống                                                 & Lấy dữ liệu danh sách sản phẩm từ server                          \\ \cline{2-4} 
                                                                                       & 3   & Hệ thống                                                 & Hiển thị danh sách sản phẩm                                                            \\ \cline{2-4} 
                                                                                       & 4   & Người dùng                                                 & Nhấn vào sản phẩm để xem chi tiết                                                                                     \\ \cline{2-4} 
                                                                                       & 5   & Hệ thống     & \begin{tabular}[c]{@{}l@{}}Truy vấn cơ sở dữ liệu để lấy ra thông tin \\ 
chi tiết của sản phẩm thông qua productId\end{tabular}             \\ \cline{2-4}
                                                                                        & 6   & Hệ thống     &\begin{tabular}[c]{@{}l@{}} Chuyển hướng đến màn hình thông tin \\
chi tiết của sản phẩm \end{tabular}   \\ \cline{2-4}
 & 7   & Hệ thống     & Hiển thị thông tin chi tiết của product   
\\ \hline
\multicolumn{1}{|l|}{\begin{tabular}[c]{@{}l@{}}\textbf{Luồng sự kiện thay}\\ \textbf{thế}\end{tabular}} & 2a & Hệ thống                            & \begin{tabular}[c]{@{}l@{}}Không thể tải danh sách sản phẩm,\\ hiện thông báo lỗi\end{tabular} \\ \hline
\multicolumn{1}{|l|}{\textbf{Xử lý ngoại lệ}}                                                   & 5a & Hệ thống                            & \begin{tabular}[c]{@{}l@{}}Không thể hiển thị chi tiết sản phẩm,\\ hiện thông báo lỗi\end{tabular}                     \\ \hline
\end{tabular}
}
\\
\multicolumn{4}{|l|}{{}}
\\ \hline
\multicolumn{1}{|l|}{\textbf{Hậu điều kiện}}  & \multicolumn{3}{l|}{Không}                                  \\ \hline
\end{tabular}
\caption{Bảng đặc tả Use case "Xem sản phẩm"}
\label{tab:viewProduct}
\end{table}
\FloatBarrier
\newpage
* Dữ liệu đầu vào:

\begin{table}[h!]
\centering
\begin{tabular}{|c|c|c|p{3.5cm}|p{1.5cm}|p{4cm}|}
\hline
\textbf{STT} & \textbf{Trường dữ liệu} & \textbf{Bắt buộc} & \textbf{Mô tả} & \textbf{Định dạng hiển thị} & \textbf{Ví dụ} \\ \hline
1 & title & Có & Tên sản phẩm & Chuỗi & Máy chạy bộ \\ \hline
2 & description & Không & Mô tả chi tiết về sản phẩm & String & Công suất: 2.0HP \\ \hline
3 & detail & Không & Mô tả một số thông tin khác như nơi sản xuất, năm sản xuất & String & Sản xuất tại Trung Quốc \\ \hline
4 & price & Có & Giá của sản phẩm đó & Number & 28,000,000 \\ \hline
5 & brandName & Có & Thương hiệu của sản phẩm & String & Samsung \\ \hline
6 & shopName & Không & Tên cửa hàng đang bán sản phẩm đó & String & LoveShop \\ \hline
7 & images & Không & Danh sách hình ảnh của sản phẩm & Danh sách & \url{https://firebasestorage.googleapis.com/v0/b/da-gr1.appspot.com/o/Banners%2Forder?alt=media&token=...} \\ \hline
\end{tabular}
\caption{Dữ liệu đầu vào cho Use case "Xem sản phẩm"}
\label{tab:viewProductInput}
\end{table}
















\FloatBarrier
\subsection{Đặc tả use case "Thêm sản phẩm vào giỏ hàng"}
\hfill
\begin{table}[ht]
\centering
\begin{tabular}{|llrl|}
\hline
\multicolumn{1}{|l|}{\textbf{Mã Use case}}    & \multicolumn{1}{l|}{2} & \multicolumn{1}{r|}{\textbf{Tên Use case}} &  \multicolumn{1}{r|}{Thêm sản phẩm vào giỏ hàng} \\ \hline

\multicolumn{1}{|l|}{\textbf{Tác nhân}}       & \multicolumn{3}{l|}{Người dùng} \\ \hline

\multicolumn{1}{|l|}{\textbf{Tiền điều kiện}} & \multicolumn{3}{l|}{%
\begin{tabular}[c]{@{}l@{}}Người dùng đang xem trang chi tiết sản phẩm \\ hoặc ở màn hình danh sách sản phẩm\end{tabular}} \\ \hline

\multicolumn{4}{|l|}{} \\ 

\multicolumn{4}{|l|}{
\begin{tabular}{|c|l|c|l|}
\hline
\multirow{5}{*}{\textbf{Luồng sự kiện chính}} & \textbf{STT} & \begin{tabular}[c]{@{}c@{}}\textbf{Thực hiện} \\ \textbf{bởi}\end{tabular} & \textbf{Hành động} \\ \cline{2-4} 
& 1 & Người dùng & \begin{tabular}[c]{@{}l@{}}Chọn số lượng sản phẩm muốn mua \\ (nếu có tùy chọn)\end{tabular} \\ \cline{2-4}
& 2 & Người dùng & Nhấn nút "Thêm vào giỏ hàng" \\ \cline{2-4}
& 3 & Hệ thống & \begin{tabular}[c]{@{}l@{}}Kiểm tra tình trạng sản phẩm \\ (còn hàng hay hết hàng)\end{tabular} \\ \cline{2-4}
& 4 & Hệ thống & \begin{tabular}[c]{@{}l@{}}Nếu còn hàng, thêm sản phẩm vào \\ giỏ hàng\end{tabular} \\ \cline{2-4}
& 5 & Hệ thống & \begin{tabular}[c]{@{}l@{}}Hiển thị thông báo \\ "Đã thêm sản phẩm vào giỏ hàng"\end{tabular} \\ \hline

\multicolumn{1}{|l|}{\begin{tabular}[c]{@{}l@{}}\textbf{Luồng sự kiện thay}\\ \textbf{thế}\end{tabular}} & 4a & Hệ thống & \begin{tabular}[c]{@{}l@{}}Sản phẩm đã hết hàng, hiển thị thông báo \\ "Sản phẩm hiện đã hết hàng"\end{tabular} \\ \hline

\multicolumn{1}{|l|}{\textbf{Xử lý ngoại lệ}} & 4b & Hệ thống & \begin{tabular}[c]{@{}l@{}}Lỗi khi thêm sản phẩm vào giỏ hàng, \\ hiển thị thông báo lỗi\end{tabular} \\ \hline
\end{tabular}
} \\

\multicolumn{4}{|l|}{} \\ \hline

\multicolumn{1}{|l|}{\textbf{Hậu điều kiện}}  & \multicolumn{3}{l|}{%
\begin{tabular}[c]{@{}l@{}}Sản phẩm được thêm vào giỏ hàng của Người mua hàng. \\ Hệ thống cập nhật số lượng sản phẩm trong giỏ hàng.\end{tabular}} \\ \hline
\end{tabular}
\caption{Bảng đặc tả Use case "Thêm sản phẩm vào giỏ hàng"}
\label{tab:addToCart}
\end{table}

\FloatBarrier
* Dữ liệu đầu vào:

\begin{table}[h!]
\centering
\begin{tabular}{|c|c|c|p{3.5cm}|p{1.5cm}|p{4cm}|}
\hline
\textbf{STT} & \textbf{Trường dữ liệu} & \textbf{Bắt buộc} & \textbf{Mô tả} & \textbf{Định dạng hiển thị} & \textbf{Ví dụ} \\ \hline
1 & userId & Có & ID của người mua hàng thực hiện thao tác thêm vào giỏ & String & user123 \\ \hline
2 & productId & Có & ID của sản phẩm cần được thêm vào giỏ hàng & String & prod789 \\ \hline
3 & quantity & Có & Số lượng sản phẩm mà người dùng muốn mua & Number & 2 \\ \hline
4 & addedAt & Không & Thời điểm sản phẩm được thêm vào giỏ hàng & DateTime & 2025-05-23T14:32:00 \\ \hline
5 & selectedOptions & Không & Tùy chọn của sản phẩm như kích cỡ, màu sắc (nếu có) & JSON & \{"size": "M", "color": "red"\} \\ \hline
\end{tabular}
\caption{Dữ liệu đầu vào cho Use case "Thêm sản phẩm vào giỏ hàng"}
\label{tab:addToCartInput}
\end{table}
\FloatBarrier

\subsection{Đặc tả use case "Đặt hàng"}
\hfill

\begin{table}[ht]
\centering
\begin{tabular}{|llrl|}
\hline
\multicolumn{1}{|l|}{\textbf{Mã Use case}}    & \multicolumn{1}{l|}{3} & \multicolumn{1}{r|}{\textbf{Tên Use case}} &  \multicolumn{1}{r|}{Đặt hàng} \\ \hline

\multicolumn{1}{|l|}{\textbf{Tác nhân}}       & \multicolumn{3}{l|}{Người dùng} \\ \hline

\multicolumn{1}{|l|}{\textbf{Tiền điều kiện}} & \multicolumn{3}{l|}{
\begin{tabular}[c]{@{}l@{}}Người dùng đã đăng nhập vào hệ thống  và giỏ hàng có ít nhất một sản phẩm\end{tabular}} \\ \hline

\multicolumn{4}{|l|}{} \\ 

\multicolumn{4}{|l|}{
\begin{tabular}{|c|l|c|l|}
\hline
\multirow{5}{*}{\textbf{Luồng sự kiện chính}} & \textbf{STT} & \begin{tabular}[c]{@{}c@{}}\textbf{Thực hiện} \\ \textbf{bởi}\end{tabular} & \textbf{Hành động} \\ \cline{2-4} 
& 1 & Người dùng & \begin{tabular}[c]{@{}l@{}}Xem giỏ hàng và click nút đặt hàng \end{tabular} \\ \cline{2-4}
& 2 & Hệ thống & Chuyển đến màn hình đặt hàng \\ \cline{2-4}
& 3 & Hệ thống & \begin{tabular}[c]{@{}l@{}}Hiển thị tóm tắt thông tin của các sản phẩm  trong \\  đơn hàng  hiển thị giá trị đơn hàng, tự động  tính \\ phí  ship,VAT và tổng giá trị sau cùng\end{tabular} \\ \cline{2-4}
& 4 & Người dùng & \begin{tabular}[c]{@{}l@{}}Chọn nút áp dụng voucher\end{tabular} \\ \cline{2-4}
& 5 & Hệ thống & \begin{tabular}[c]{@{}l@{}}Tính toán và đưa ra danh sách các voucher  mà\\  người đó đủ điều kiện áp dụng\end{tabular} \\ \cline{2-4}
& 6 & Người dùng & \begin{tabular}[c]{@{}l@{}}Chọn 1 loại voucher bất kì trong danh sách  \end{tabular} \\ \cline{2-4}
& 7 & Hệ thống & \begin{tabular}[c]{@{}l@{}}Tính toán và cập nhật lại giá trị đơn hàng, thêm \\ vào list voucher được áp dụng và  hiển thị list lên\\ giao diện \\ \end{tabular} \\ \cline{2-4}
& 8 & Người dùng & \begin{tabular}[c]{@{}l@{}}Chọn phương thức thanh toán  \end{tabular} \\ \cline{2-4}
& 9 & Người dùng & \begin{tabular}[c]{@{}l@{}}Chọn địa chỉ giao hàng \end{tabular} \\ \cline{2-4}
& 10 & Hệ thống & \begin{tabular}[c]{@{}l@{}}Hiển thị đầy đủ thông tin phương thức thanh toán\\, địa chỉ giao hàng lên giao diện \\ 
\end{tabular} \\ \cline{2-4}
& 11 & Người dùng & \begin{tabular}[c]{@{}l@{}}Chọn nút thanh toán \end{tabular} \\ \cline{2-4}
& 12 & Hệ thống & \begin{tabular}[c]{@{}l@{}}Kiểm tra thông tin người dùng đã hợp lệ chưa \\ chuyển sang màn hình thanh toán \\ 
\end{tabular} \\ 
\hline

\multicolumn{1}{|l|}{\begin{tabular}[c]{@{}l@{}}\textbf{Luồng sự kiện thay}\\ \textbf{thế}\end{tabular}}
& 5a & Hệ thống & \begin{tabular}[c]{@{}l@{}}Nếu không có voucher nào  có thể áp dụng,\\  hiển thị thông báo "Không thể áp dụng voucher"\end{tabular} \\ \cline{2-4}
& 7a & Hệ thống & \begin{tabular}[c]{@{}l@{}}Nếu 1 voucher sau khi áp dụng, giá trị đơn hàng \\  đã về 0 mà người đó   vẫn tiếp tục áp dụng thì đưa\\ ra  cảnh báo  "Không cần áp dụng thêm voucher" \end{tabular} \\ \cline{2-4}
& 10a & Hệ thống & \begin{tabular}[c]{@{}l@{}}Nếu người dùng chưa chọn thông tin địa chỉ\\ hoặc phương thức thanh toán,  hiển thị thông báo \\ "Bạn cần chọn đầy đủ thông tin"\end{tabular} \\ \cline{2-4}
& 12a & Hệ thống & \begin{tabular}[c]{@{}l@{}}Nếu người dùng không chọn nút thanh toán,\\ mà back về trang trước thì hệ thống cũng xóa các\\ thông tin đã chọn và điều hướng về màn hình đó\end{tabular} \\ 

\hline

\multicolumn{1}{|l|}{\textbf{Xử lý ngoại lệ}} & 9a & Hệ thống & \begin{tabular}[c]{@{}l@{}}Nếu chưa thêm địa chỉ nào, hiển thị thông báo lỗi \\ và yêu cầu người dùng thêm một địa chỉ mới\end{tabular} \\ \hline
\end{tabular}
} \\

\multicolumn{4}{|l|}{} \\ \hline

\multicolumn{1}{|l|}{\textbf{Hậu điều kiện}}  & \multicolumn{3}{l|}{%
\begin{tabular}[c]{@{}l@{}}Không\end{tabular}} \\ \hline
\end{tabular}

\caption{Bảng đặc tả Use case "Đặt hàng"}
\label{tab:addToCart}
\end{table}





\FloatBarrier
\subsection{Đặc tả use case "Thanh toán điện tử"}
\hfill

\begin{table}[ht]
\centering
\begin{tabular}{|llrl|}
\hline
\multicolumn{1}{|l|}{\textbf{Mã Use case}}    & \multicolumn{1}{l|}{4} & \multicolumn{1}{r|}{\textbf{Tên Use case}} &  \multicolumn{1}{r|}{Thanh toán điện tử} \\ \hline

\multicolumn{1}{|l|}{\textbf{Tác nhân}}       & \multicolumn{3}{l|}{Người dùng}                             \\ \hline
\multicolumn{1}{|l|}{\textbf{Tiền điều kiện}} & \multicolumn{3}{l|}{\begin{tabular}[c]{@{}l@{}}Người dùng đã đăng nhập vào hệ thống\\ Người dùng đã đặt hàng thành công\end{tabular}}                                \\ \hline
\multicolumn{4}{|l|}{}
                    \\ 
\multicolumn{4}{|l|}{
\begin{tabular}{|c|l|c|l|}
\hline
\multirow{8}{*}{\textbf{Luồng sự kiện chính}}                                                   & \textbf{STT} & \begin{tabular}[c]{@{}c@{}}\textbf{Thực hiện} \\ \textbf{bởi}\end{tabular} & \textbf{Hành động}                                                                                                 \\ \cline{2-4} 
 & 1   & Ngươi dùng  & \begin{tabular}[c]{@{}l@{}}Chọn nút thanh toán\end{tabular}             \\ \cline{2-4} 
 & 2   & Hệ thống                                                 & \begin{tabular}[c]{@{}l@{}}Chuyển đến màn hình thanh toán \end{tabular}                           \\ \cline{2-4} 
 & 3   & Hệ thống                                                 & \begin{tabular}[c]{@{}l@{}}Gọi đến API thanh toán Stripe và hiển thị\\popup thanh toán \end{tabular}           \\ \cline{2-4}                                                          
& 4   & Người dùng                                                &   \begin{tabular}[c]{@{}l@{}}Nhập đầy đủ các thông tin như số thẻ, thời\\ hạn, CVC, địa chỉ\end{tabular}                      \\ \cline{2-4} 
& 5   & Người dùng                                                &   \begin{tabular}[c]{@{}l@{}}Kiểm tra lại thông tin đã nhập, xem số tiền\\phải trả và bấm thanh toán ngay\end{tabular} \\ \cline{2-4} 
& 6   & Hệ thống                                                &   \begin{tabular}[c]{@{}l@{}}Gọi API, tạo giao dịch mới và gửi thông tin\\ giao dịch sang cho Stripe xử lý\end{tabular} \\ \cline{2-4} 
& 7   & Hệ thống                                                &   \begin{tabular}[c]{@{}l@{}}Nhận phản hồi từ Stripe, hiển thị kết quả \\ lên giao diện\end{tabular} \\ \cline{2-4} 
& 8   & Hệ thống                                                &   \begin{tabular}[c]{@{}l@{}}Gửi thông báo đơn hàng đã thanh toán thành \\ công và điều hướng về màn hình home\end{tabular}
\\ \hline
\multicolumn{1}{|l|}{\begin{tabular}[c]{@{}l@{}}\textbf{Luồng sự kiện thay}\\ \textbf{thế}\end{tabular}}
& 6a & Hệ thống & \begin{tabular}[c]{@{}l@{}}Nếu người dùng nhập thiếu thông tin, hiển \\thị cảnh báo và yêu cầu người dùng phải\\ nhập đủ \end{tabular} \\ \cline{2-4}
& 7a & Hệ thống & \begin{tabular}[c]{@{}l@{}}Nếu có lỗi không thực hiện được giao dịch\\ thì hiện thông báo lỗi, ẩn popup và\\ trở về màn hình trước đó \end{tabular} \\
\hline
\multicolumn{1}{|l|}{\textbf{Xử lý ngoại lệ}}                                                   & 3a & \multicolumn{1}{c|}{Hệ thống}                            & \begin{tabular}[c]{@{}l@{}}Không gọi đc API Stripe, hiện thông báo lỗi\end{tabular}                     \\ \hline
\end{tabular}
}
\\
\multicolumn{4}{|l|}{{}}
\\ \hline
\multicolumn{1}{|l|}{\textbf{Hậu điều kiện}}  & \multicolumn{3}{l|}{Không}                                  \\ \hline
\end{tabular}
\caption{Bảng đặc tả Use case "Thanh toán điện tử"}
\label{tab:payment}
\end{table}
\FloatBarrier
* Dữ liệu đầu vào:

\begin{table}[ht]
\centering
\begin{tabular}{|l|l|l|l|l|l|}
\hline
\multicolumn{1}{|c|}{STT} & \multicolumn{1}{c|}{\begin{tabular}[c]{@{}c@{}}Trường dữ\\ liệu\end{tabular}} & \multicolumn{1}{c|}{Bắt buộc} & \multicolumn{1}{c|}{Mô tả}                                                    & \multicolumn{1}{c|}{\begin{tabular}[c]{@{}c@{}}Định dạng\\ hiển thị\end{tabular}} & \multicolumn{1}{c|}{Ví dụ}                                     \\ \hline
1                         & Số thẻ                                                                     & Có                            & \begin{tabular}[c]{@{}l@{}}Số thẻ ngân hàng\\ liên kết \end{tabular}                                                             & String                                                                            &  \begin{tabular}[c]{@{}l@{}} 4242 4242 \\ 4242 4242 \end{tabular}                                                            \\ \hline
2                         &                 Hiệu lực                                                    & Có                            & \begin{tabular}[c]{@{}l@{}}Thời gian còn hiệu \\lực của thẻ \end{tabular}      & Date                                                                            & \begin{tabular}[c]{@{}l@{}}12/27\end{tabular}                                                       \\ \hline
3                         &  CVC                                                                & Có                         & \begin{tabular}[c]{@{}l@{}}Mật mã bí mật\\của thẻ \end{tabular} & Number                                                                              & 222                                                     \\ \hline
4                         & Địa chỉ                                                                      & Có                        & \begin{tabular}[c]{@{}l@{}}Quốc gia hoặc khu vực\\ sử dụng\end{tabular}            & String                                                                            & \begin{tabular}[c]{@{}l@{}}Việt Nam\end{tabular}                                                         \\ \hline
\end{tabular}
\caption{Dữ liệu đầu vào cho Use case "Thanh toán điện tử"}
\label{tab:paymentInput}
\end{table}
\FloatBarrier

\subsection{Đặc tả use case "Tham gia nhiệm vụ"}
\hfill

\begin{table}[h!]
\centering
\begin{tabular}{|llrl|}
\hline
\multicolumn{1}{|l|}{\textbf{Mã Use case}}    & \multicolumn{1}{l|}{5} & \multicolumn{1}{r|}{\textbf{Tên Use case}} &  \multicolumn{1}{r|}{Tham gia nhiệm vụ} \\ \hline

\multicolumn{1}{|l|}{\textbf{Tác nhân}}       & \multicolumn{3}{l|}{Người dùng}                             \\ \hline
\multicolumn{1}{|l|}{\textbf{Tiền điều kiện}} & \multicolumn{3}{l|}{\begin{tabular}[c]{@{}l@{}}Người dùng đã đăng nhập vào hệ thống \end{tabular}}                                \\ \hline
\multicolumn{4}{|l|}{}
                    \\ 
\multicolumn{4}{|l|}{
\begin{tabular}{|c|l|c|l|}
\hline
\multirow{7}{*}{\textbf{Luồng sự kiện chính}}                                                   & \textbf{STT} & \begin{tabular}[c]{@{}c@{}}\textbf{Thực hiện} \\ \textbf{bởi}\end{tabular} & \textbf{Hành động}                                                                                                 \\ \cline{2-4} 
 & 1   & Người dùng    & \begin{tabular}[c]{@{}l@{}}Chọn màn hình danh sách nhiệm vụ \end{tabular}             \\ \cline{2-4} 
 & 2   & Hệ thống  & \begin{tabular}[c]{@{}l@{}}Truy vấn server để lấy danh sách nhiệm vụ\\  và trạng thái thực hiện nhiệm vụ của người\\dùng đó\end{tabular}             \\ \cline{2-4} 
 & 3   & Người dùng                                                 & \begin{tabular}[c]{@{}l@{}}Bấm vào 1 nhiệm vụ để xem thông tin chi tiết \\ nhiệm vụ\end{tabular}                           \\ \cline{2-4} 
 & 4   & Hệ thống                                                 &\begin{tabular}[c]{@{}l@{}}Lấy id của nhiệm vụ gửi lên server và\\ lấy thông tin \\ chi tiết của nhiệm vụ, điều hướng đến màn\\ hình chi tiết\end{tabular}                                                                                             \\ \cline{2-4} 
 & 5   & Người dùng                                                 & \begin{tabular}[c]{@{}l@{}}Xem thông tin chi tiết gồm tên, trạng thái \\ nhiệm vụ, điểm nhận được và deadline \\ hoàn thành\end{tabular}                                                                                     \\ \cline{2-4} 
 & 6   & Người dùng                                                 & \begin{tabular}[c]{@{}l@{}}Bấm nút bắt đầu nhiệm vụ \end{tabular}   \\ \cline{2-4}    & 7   & Hệ thống                                                 & \begin{tabular}[c]{@{}l@{}}Hiển thị đồng hồ đếm ngược của nhiệm vụ,\\ cập nhật trạng thái của nhiệm vụ, lưu thông\\ tin nhiệm vụ của người dùng lên server\end{tabular}   \\ \cline{2-4}                                                       & 8   & Hệ thống                                                 & \begin{tabular}[c]{@{}l@{}}Hệ thống tạo tracking để theo dõi việc thực\\ hiện nhiệm vụ của người đó, khi người dùng\\ hoàn thành 1 step thì giao diện cũng sẽ tự \\động cập nhật và lưu lên server\end{tabular}   \\ \cline{2-4}       & 9   & Hệ thống                                                 & \begin{tabular}[c]{@{}l@{}}Hệ thống theo dõi và thấy nếu người dùng\\ hoàn  thành nhiệm vụ trước khi hết thời \\ gian đếm ngược thì sẽ tự động gửi thông báo \\ nhiệm\\ vụ  thành công và cộng điểm cho người \\ dùng, lưu lên server\end{tabular} \\ \hline                        
                                                                                                                                                                              
\multirow{2}{*}{\begin{tabular}[c]{@{}l@{}}\textbf{Luồng sự kiện thay}\\ \textbf{thế}\end{tabular}} & 7a & \multicolumn{1}{c|}{Hệ thống}                            & \begin{tabular}[c]{@{}l@{}}Nếu 1 nhiệm vụ đã bắt đầu rồi thì nó sẽ \\ hiển thị thời gian còn lại\end{tabular} \\ \cline{2-4} & 9a & \multicolumn{1}{c|}{Hệ thống}                            & \begin{tabular}[c]{@{}l@{}}Nếu hết thời gian mà nhiệm vụ chưa hoàn\\ thành thì nó tự động cập nhật trạng thái, \\ẩn đồng hồ đếm ngược và lưu lên server\end{tabular} \\ \hline
\multicolumn{1}{|l|}{\textbf{Xử lý ngoại lệ}}                                                   & 2a & \multicolumn{1}{c|}{Hệ thống}                            & \begin{tabular}[c]{@{}l@{}}Không lấy được dữ liệu,hiện thông báo lỗi\end{tabular}                     \\ \hline
\end{tabular}
}
\\
\multicolumn{4}{|l|}{{}}
\\ \hline
\multicolumn{1}{|l|}{\textbf{Hậu điều kiện}}  & \multicolumn{3}{l|}{Không}                                  \\ \hline
\end{tabular}
\caption{Bảng đặc tả Use case Tham gia nhiệm vụ}
\label{tab:joinMission}
\end{table}

\FloatBarrier

\subsection{Đặc tả use case "Mời bạn bè tham gia nhóm"}
\hfill

\begin{table}[h!]
\centering
\begin{tabular}{|llrl|}
\hline
\multicolumn{1}{|l|}{\textbf{Mã Use case}}    & \multicolumn{1}{l|}{6} & \multicolumn{1}{r|}{\textbf{Tên Use case}} &  \multicolumn{1}{r|}{Mời bạn bè tham gia nhóm} \\ \hline

\multicolumn{1}{|l|}{\textbf{Tác nhân}}       & \multicolumn{3}{l|}{Người dùng}                             \\ \hline
\multicolumn{1}{|l|}{\textbf{Tiền điều kiện}} & \multicolumn{3}{l|}{\begin{tabular}[c]{@{}l@{}}Người dùng đã đăng nhập vào hệ thống\\ Người dùng đã tạo 1 nhóm săn sale\end{tabular}}                                \\ \hline
\multicolumn{4}{|l|}{}
                    \\ 
\multicolumn{4}{|l|}{
\begin{tabular}{|c|l|c|l|}
\hline
\multirow{9}{*}{\textbf{Luồng sự kiện chính}}                                                   & \textbf{STT} & \begin{tabular}[c]{@{}c@{}}\textbf{Thực hiện} \\ \textbf{bởi}\end{tabular} & \textbf{Hành động}                                                                                                 \\ \cline{2-4} 
 & 1   & Người dùng    & \begin{tabular}[c]{@{}l@{}}Chọn nút mời bạn bè tham gia nhóm\end{tabular}             \\ \cline{2-4}                                    & 2   & Hệ thống  & \begin{tabular}[c]{@{}l@{}}Điều hướng đến màn hình mời, hiển thị \\ danh sách bạn bè kèm nút invite \end{tabular}             \\ \cline{2-4}
 & 3   & Người dùng                                                 & \begin{tabular}[c]{@{}l@{}}Chọn nút invite và mời một người bạn \\ vào nhóm\end{tabular}                           \\ \cline{2-4}                   
& 4   & Hệ thống                                                 & \begin{tabular}[c]{@{}l@{}}Gửi request đến người bạn đó\\ thông báo đã gửi yêu cầu thành công\end{tabular}                                                                                             \\ \cline{2-4}                                                                   & 5   & Hệ thống                                                 & \begin{tabular}[c]{@{}l@{}}Khi có phản hồi từ bạn bè,\\ gửi thông báo đã chấp nhận hoặc đã \\từ chối lời mời\end{tabular} \\ \cline{2-4}                                                                                                            
   & 6   & Hệ thống     & \begin{tabular}[c]{@{}l@{}}Cập nhật lại trạng thái, số thành viên \\  của nhóm hiện tại\end{tabular} \\ \cline{2-4}                    & 7   & Hệ thống     & \begin{tabular}[c]{@{}l@{}}Nếu đã đủ thành viên, gửi thông báo \\  săn sale thành công đến tất cả các \\  thành viên trong nhóm, tự động thêm \\  voucher săn sale vào kho của mỗi \\  thành viên\end{tabular}   \\ \hline                                                                                                                                                                              
                                                                                                                                                                                                                                              
\multirow{2}{*}{\begin{tabular}[c]{@{}l@{}}\textbf{Luồng sự kiện thay}\\ \textbf{thế}\end{tabular}}
& 2a & \multicolumn{1}{c|}{Hệ thống}                            & \begin{tabular}[c]{@{}l@{}}Nếu chưa có người bạn nào sẽ thông \\  báo và có nút để điều hướng đến\\ màn hình kết bạn\end{tabular}
\\ \hline
\multicolumn{1}{|l|}{\textbf{Xử lý ngoại lệ}}                                                   & 4a & \multicolumn{1}{c|}{Hệ thống}                            & \begin{tabular}[c]{@{}l@{}}Không gửi được yêu cầu, hiện thông \\ báo lỗi\end{tabular}                     \\ \hline
\end{tabular}
}
\\
\multicolumn{4}{|l|}{{}}
\\ \hline
\multicolumn{1}{|l|}{\textbf{Hậu điều kiện}}  & \multicolumn{3}{l|}{Không}                                  \\ \hline
\end{tabular}
\caption{Bảng đặc tả Use case "Mời bạn bè tham gia nhóm"}
\label{tab:updatePlay}
\end{table}
\FloatBarrier



\subsection{Đặc tả use case "Viết bình luận"}
\hfill
\begin{table}[ht]
\centering
\begin{tabular}{|llrl|}
\hline
\multicolumn{1}{|l|}{\textbf{Mã Use case}}    & \multicolumn{1}{l|}{7} & \multicolumn{1}{r|}{\textbf{Tên Use case}} &  \multicolumn{1}{r|}{Viết bình luận} \\ \hline

\multicolumn{1}{|l|}{\textbf{Tác nhân}}       & \multicolumn{3}{l|}{Người dùng}                             \\ \hline
\multicolumn{1}{|l|}{\textbf{Tiền điều kiện}} & \multicolumn{3}{l|}{\begin{tabular}[c]{@{}l@{}}Người dùng đăng nhập vào hệ thống  \end{tabular}}                                \\ \hline
\multicolumn{4}{|l|}{}
                    \\ 
\multicolumn{4}{|l|}{
\begin{tabular}{|c|l|c|l|}
\hline
\multirow{9}{*}{\textbf{Luồng sự kiện chính}}                                                   & \textbf{STT} & \begin{tabular}[c]{@{}c@{}}\textbf{Thực hiện} \\ \textbf{bởi}\end{tabular} & \textbf{Hành động}                                                                                                 \\ \cline{2-4} 
 & 1   & Người dùng    & \begin{tabular}[c]{@{}l@{}}Chọn nút viết bình luận\end{tabular}             \\ \cline{2-4}                                    & 2   & Hệ thống  & \begin{tabular}[c]{@{}l@{}}Điều hướng đến màn hình review 1 sản \\ phẩm \end{tabular}             \\ \cline{2-4}
 & 3   & Hệ thống                                                & \begin{tabular}[c]{@{}l@{}}Hiển thị form review sản phẩm gồm tên \\ và ảnh của sản phẩm\end{tabular}                           \\ \cline{2-4}                   
& 4   & Người dùng                                                 & \begin{tabular}[c]{@{}l@{}}Chọn số sao vote cho sản phẩm \end{tabular}                                                                                             \\ \cline{2-4}                                                                   & 5   & Người dùng                                                 & \begin{tabular}[c]{@{}l@{}}Nhập bình luận về sản phẩm\end{tabular} \\ \cline{2-4}                                                                                                            
   & 6   & Hệ thống     & \begin{tabular}[c]{@{}l@{}}Kiểm soát độ dài nội dung bình luận,\\hiển thị thông tin người dùng đã\\ nhập\end{tabular} \\ \cline{2-4} 
   & 7   & Người dùng     & \begin{tabular}[c]{@{}l@{}}Chọn nút upload ảnh hoặc video \end{tabular}\\ \cline{2-4} 
 & 8  & Hệ thống     & \begin{tabular}[c]{@{}l@{}}Hiển thị popup để người dùng  lựa chọn\\ up ảnh hay up video  up từ \\ thư viện hay chụp mới \end{tabular}\\ \cline{2-4}   
 &  9 & Hệ thống     & \begin{tabular}[c]{@{}l@{}}Điều hướng đến thư viện ảnh hoặc \\ camera, lưu data người dùng chọn \\ và quay lại form nhập bình luận \\ thông báo đã up ảnh, video thành công \end{tabular}\\ \cline{2-4}   
  &  10 & Người dùng     & \begin{tabular}[c]{@{}l@{}}Chọn chế độ up bình luận là công khai \\ hay ẩn danh\end{tabular}\\ \cline{2-4}  
  &  11 & Hệ thống     & \begin{tabular}[c]{@{}l@{}}Kiểm tra sự hợp lệ của form gửi dữ\\ liệu và lưu lên server và quay\\ lại màn hình xem bình luận \\ về sản phẩm\end{tabular}
    
   \\ \hline
                                                                                                
\multicolumn{1}{|l|}{\begin{tabular}[c]{@{}l@{}}\textbf{Luồng sự kiện thay}\\ \textbf{thế}\end{tabular}} & 6a & \multicolumn{1}{c|}{Hệ thống}                            & \begin{tabular}[c]{@{}l@{}}Nếu người dùng nhập bình luận quá dài\\sẽ thông báo và không cho nhập thêm \end{tabular} \\ \hline
\multicolumn{1}{|l|}{\textbf{Xử lý ngoại lệ}}                                                   & 11a & \multicolumn{1}{c|}{Hệ thống}                            & \begin{tabular}[c]{@{}l@{}}Không lưu được dữ liệu, hiện thông báo lỗi\end{tabular}                     \\ \hline
\end{tabular}
}
\\
\multicolumn{4}{|l|}{{}}
\\ \hline
\multicolumn{1}{|l|}{\textbf{Hậu điều kiện}}  & \multicolumn{3}{l|}{Không}                                  \\ \hline
\end{tabular}
\caption{Bảng đặc tả Use case "Viết bình luận"}
\label{tab:writeReview} 
\end{table}
\FloatBarrier
* Dữ liệu đầu vào:

\begin{table}[h!]
\centering
\renewcommand{\arraystretch}{1.4}
\resizebox{\textwidth}{!}{%
\begin{tabular}{|c|p{2.5cm}|c|p{4.5cm}|p{3cm}|p{4.5cm}|}
\hline
\textbf{STT} & \textbf{Trường dữ liệu} & \textbf{Bắt buộc} & \textbf{Mô tả} & \textbf{Định dạng hiển thị} & \textbf{Ví dụ} \\ \hline
1 & rating & Không & Số sao đánh giá cho sản phẩm & Number & 5 \\ \hline
2 & review & Không & Nội dung bình luận & String & Sản phẩm tuyệt lắm \\ \hline
3 & image & Không & Các ảnh người dùng upload & Img & \url{https://url/img} \\ \hline
4 & video & Không & Các video người dùng upload & Video & \url{https://www.youtube.com/watch?v=5qt3enCkcwo&list=RDMMSASAH-gigmQ&index=27} \\ \hline
\end{tabular}%
}
\caption{Dữ liệu đầu vào cho Use case "Viết bình luận"}
\label{tab:writeReviewInput}
\end{table}
\FloatBarrier


\section{Yêu cầu phi chức năng}
\label{section:2.4}
% Trong phần này, sinh viên đưa ra các yêu cầu khác nếu có, bao gồm các yêu cầu phi chức năng như hiệu năng, độ tin cậy, tính dễ dùng, tính dễ bảo trì, hoặc các yêu cầu về mặt kỹ thuật như về CSDL, công nghệ sử dụng, v.v.
Ngoài các chức năng đã được liệt kê ở phần trước, hệ thống cũng cần phải thêm một số yêu cầu phi chức năng để đáp ứng hơn nữa nhu cầu sử dụng của người dùng:
\subsection{Yêu cầu về giao diện và trải nghiệm người dùng}
Giao diện của ứng dụng cần phải được thiết kế bắt mắt, bố cục phù hợp, dễ sử dụng, màu sắc hài hòa, dễ nhìn và các nút bấm cần đủ to để người dùng có thể nhấn một cách dễ dàng trên thiết bị di động. Ngoài ra, khi người dùng gặp lỗi khi nhập liệu hay gửi yêu cầu, sẽ luôn có các thông báo lỗi được hiển thị để giúp người dùng biết được nội dung lỗi và chỉnh sửa lại sao cho phù hợp.

\subsection{Yêu cầu về tính dễ bảo trì}
Ứng dụng cần phải có tính dễ bảo trì để dễ dàng chỉnh sửa các thành phần riêng biệt mà không làm ảnh hưởng đến cả hệ thống. Phần client và server được phân tách riêng biệt với nhau bằng các thư mục riêng biệt. Ngoài ra phần giao diện được phân chia thành các component giúp dễ dành trong việc quản lý code và chỉnh sửa sau này.

\subsection{Yêu cầu về bảo mật}
Ứng dụng cần đảm bảo tính bảo mật cao để bảo vệ thông tin cá nhân của người dùng cũng như dữ liệu liên quan đến đơn hàng, sản phẩm và thanh toán. Firebase cung cấp nhiều cơ chế bảo mật để hỗ trợ điều này:
\begin{itemize}
     

\item Xác thực người dùng (Firebase Authentication \cite{FirebaseAuth}): Tất cả người dùng phải đăng nhập thông qua các phương thức xác thực an toàn như Email/Password, Google, hoặc OTP qua điện thoại. Sau khi đăng nhập, người dùng sẽ được gán một mã định danh duy nhất (UID), giúp xác định danh tính trong các truy vấn dữ liệu.

\item Quy tắc bảo mật cho Firestore \cite{FirebaseFirestore} (Security Rules): Ứng dụng sử dụng Firestore Security Rules để kiểm soát quyền đọc/ghi dữ liệu theo UID. Chỉ người dùng hợp lệ mới có thể truy cập dữ liệu của chính họ, ví dụ: đơn hàng, giỏ hàng, địa chỉ giao hàng,…

\item Realtime Database / Firestore \cite{FirebaseFirestore} đã mã hóa dữ liệu: Firebase tự động mã hóa dữ liệu ở trạng thái lưu trữ (at rest) và khi truyền tải (in transit), đảm bảo an toàn khi giao tiếp qua mạng. 
\end{itemize}
\subsection{Yêu cầu về hiệu năng}
Ứng dụng cần phải hoạt động mượt mà, ổn định ít giật lag, thời gian chờ cho mỗi tác vụ gửi request và nhận response giữa client và server phải đủ nhanh kể cả khi lượng dữ liệu truyền đi hoặc nhận về lớn.

Như vậy, ở trong chương 2 này em đã đưa ra các chức năng của hệ thống thông qua biểu đồ use case với 2 tác nhân chính là người dùng và quản trị viên. Ngoài ra, em cũng đã nêu ra một số yêu cầu phi chức năng cần có của hệ thống. Tiếp đến, em sẽ giới thiệu một số công nghệ được sử dụng để thực hiện các chức năng và các yêu cầu phi chức năng trên.


%%%%%%%%%%%%%%%%%%%%%%%%%%%%%%%%%%%

\end{document}